\begin{claim}{D.3.3}{Proved}{Every nonzero finite primary addition reaches the upper endpoint}
Fix a prime $p$. Let $V$ be a finite connected pointed complex with
nonzero reduced integral homology, all of it $p$-primary torsion.
Then
\[
 \operatorname{Rev}(B\ZZ/p\vee V)
       =W_{\ZZ[1/p]}=R(\operatorname{loc}(\mathbb S/p)).
 \tag{D.3.3.1}
\]
Consequently a proper intermediate localizer in the projective-space
interval cannot make any nonzero finite $p$-primary complex aspherical.
In particular adding a finite complex of type two, or of any higher
finite chromatic type, cannot produce such an intermediate.
\end{claim}
\begin{proof}
Write $W=\operatorname{Rev}(B\ZZ/p\vee V)$. Its two generators
are $\ZZ[1/p]$-acyclic, so $W\subseteq W_{\ZZ[1/p]}$.
Replacing $V$ by a sufficiently high suspension does not change
$W$, by reversibility. We may therefore suppose that $V$ is a
simply connected suspension. The nonzero finite spectrum
$\Sigma^\infty V$ has some finite $p$-type $n+1\geq1$.
Let $d$ be the lowest degree with $\pi_dV\ne0$, and choose an
odd integer $m\geq3$ with $m-1\geq d$.

Distinguish the ordinary nullification $P_V$ from the Goodwillie
stages $G_s=P_s^{\mathrm{Id}}(S^m_{(p)})$.
The precise periodicity input is
\cite[Theorem~4.1]{EH5AroneMahowald1999}: the natural map
\[
                 S^m_{(p)}\longrightarrow G_{p^n}
 \tag{D.3.3.2}
\]
is a $v_j$-periodic equivalence for every $0\leq j\leq n$,
where $j=0$ means a rational homotopy equivalence.
The nullification comparison is
\cite[Corollary~4.8]{EH5Bousfield2001}, in the arbitrary finite
suspension form explicitly stated in
\cite[Theorem~3.1]{EH5BarthelHeutsMeier2021}: a map of
$d$-connected $p$-local spaces is a $P_V$-equivalence if and
only if it is a $v_j$-equivalence for $0\leq j\leq n$.
All the spaces $G_s$ are $(m-1)$-connected. To check the bound,
the derivative formula in
\cite[Theorem~2.1]{EH5AroneMahowald1999} gives the spectrum of
the $s$th layer as
\[
 E_s=\bigl(\operatorname{Map}_*(\Sigma K_s,
          \Sigma^\infty(S^m_{(p)})^{\wedge s})\bigr)_{h\Sigma_s}.
\]
The finite complex $\Sigma K_s$ has dimension $s-1$, so $E_s$
has no homotopy below $sm-s+1\geq m$. Homotopy orbits preserve
this lower bound. Starting with $G_1=QS^m_{(p)}$ and using
the fiber sequences $\Omega^\infty E_s\to G_s\to G_{s-1}$
proves the assertion by induction. These spaces are $p$-local
as well, by the same induction. Thus (D.3.3.2) is a
$P_V$-equivalence and hence belongs to $W$.

We next remove the finite Goodwillie tower by $B\ZZ/p$
nullification. For $s>1$, the spectrum $E_s$ is rationally
zero, by \cite[Proposition~3.1]{EH5AroneMahowald1999}.
Since it is $p$-local, $E_s[1/p]=0$; and the preceding
connectivity estimate makes it a connected spectrum.
\cite[Theorem~3.2]{EH5Levy2025} therefore gives
\[
                    P_{B\ZZ/p}\Omega^\infty E_s\simeq *.
 \tag{D.3.3.3}
\]
Every projection $G_s\to G_{s-1}$ is consequently a
$P_{B\ZZ/p}$-equivalence. Here only the following elementary
nullification property is used: a fibration with connected
$A$-acyclic fibers is a $P_A$-equivalence. Indeed, mapping its
total space into an $A$-null target is the homotopy limit over
the base of the fiber mapping spaces; the natural maps from
the target to those mapping spaces are equivalences. This
identifies the resulting homotopy limit with the mapping space
from the base. No preservation of periodic homotopy by
$P_{B\ZZ/p}$ is asserted.

The map
\[
                 QS^m_{(p)}\longrightarrow K(\QQ,m)
 \tag{D.3.3.4}
\]
is also a $P_{B\ZZ/p}$-equivalence. Its fiber is the infinite
loop space of
\[
     F=\operatorname{fib}(\Sigma^m\mathbb S_{(p)}
                            \longrightarrow\Sigma^mH\QQ).
\]
The spectrum $F$ satisfies $F[1/p]=0$, and its lowest possible
nonzero homotopy is in degree $m-1\geq2$. Apply Levy's theorem
again. The target $K(\QQ,m)$ is $B\ZZ/p$-null, as is also
immediate from the rational cohomology of $B\ZZ/p$.

Both ordinary nullification-equivalence classes used here are
contained in $W$, since $W$ contains their generating maps.
Combining (D.3.3.2)--(D.3.3.4), we have a natural $W$-equivalence
\[
                         r:S^m_{(p)}\longrightarrow K(\QQ,m).
\]
Naturality for a map $p:S^m_{(p)}\to S^m_{(p)}$ of degree $p$
gives $r p\simeq p r$, with the lower $p$ the multiplication
automorphism of $K(\QQ,m)$. Two out of three therefore puts
the degree-$p$ map of $S^m_{(p)}$ in $W$. Its homotopy cofiber
is the usual finite Moore space $M(\ZZ/p,m)$: the comparison
from the integral degree-$p$ cofiber is an integral homology
equivalence of simply connected spaces. The reversible cofiber
criterion makes this Moore space $W$-aspherical.

Suspension and desuspension now put $M(\ZZ/p,2)$ in the
aspherical class. Its ordinary nullification kills every
simply connected space with $p$-primary homotopy groups, by
\cite[Theorem~7.5]{EH5Bousfield1997}. The Serre-class theorem
identifies these with the simply connected spaces having
$p$-primary reduced integral homology. If $X$ is any connected
space with such homology, its suspension is of this form;
reversibility makes $X$ aspherical as well. These are exactly
the aspherical spaces of $W_{\ZZ[1/p]}$. Applying the cofiber
criterion gives $W_{\ZZ[1/p]}\subseteq W$ and proves (D.3.3.1).
\end{proof}
